\section{Introduction}

GPU-centric computing now underpins both large-scale AI training and modern HPC simulation, and the corresponding data footprint has expanded from gigabyte-scale epochs to multi-terabyte and even petabyte-scale datasets and checkpoints. In this regime, storage and input pipelines frequently dominate end-to-end time-to-solution, even when accelerator compute is abundant. Empirical evidence from ML data-processing frameworks and distributed caching systems shows that data ingestion, shuffling, and staging can directly throttle GPU utilization if storage cannot sustain required concurrency and locality \cite{murray2021tfdata,hoard2018}. Similar pressures appear in HPC/AI production storage environments, where persistence and movement of large model states have become first-order system concerns \cite{zhang2025hpcstoragesurvey,rajbhandari2021zeroinfinity}.

The core architectural mismatch is that mainstream storage stacks were designed for CPU-centric orchestration. In traditional paths, the host CPU remains responsible for request submission, metadata handling, synchronization, and completion processing; data frequently traverses host-managed paths before reaching accelerator memory. Consequently, \emph{CPU remains on the critical path of I/O}, creating control-plane and software-path overheads that scale poorly with massively parallel GPU execution \cite{linux_io_uring,spdk_overview,nvme_spec}. For fine-grained or metadata-heavy access patterns, this mismatch manifests as low accelerator occupancy despite available storage bandwidth \cite{silberstein2013gpufs,shahar2016gpufs}.

Recent technologies have improved this situation but only partially. RDMA transports, GPUDirect RDMA, and GPUDirect Storage reduce copies and bypass host bounce buffers, yielding substantial gains in transfer efficiency \cite{infiniband_spec,nvidia_gdrdma,nvidia_gds_blog,nvidia_gds_overview}. However, most of these advances primarily optimize \emph{data movement} (where bytes flow and how fast), while leaving higher-level orchestration, metadata semantics, and control placement relatively unchanged. This is precisely why existing data-movement-centric reviews are insufficient for explaining end-to-end behavior in GPU-dominated data systems \cite{traub2023compstoragesurvey,zhang2025hpcstoragesurvey}.

This survey adopts a different perspective: the critical transition is from GPUs as passive data consumers to GPUs as active participants in storage interaction and, increasingly, storage control. We use \emph{GPU-integrated storage} to denote architectures where storage abstractions, control decisions, and data paths are co-designed around accelerator execution contexts rather than retrofitted onto CPU-first stacks. Systems such as GPUfs, BaM, and GeminiFS illustrate this shift by exposing GPU-visible storage interfaces or by pushing portions of I/O decision making closer to accelerator-driven workflows \cite{silberstein2013gpufs,qureshi2023bam,geminifs2025}.

Accordingly, the scope of this survey is \emph{not} limited to copy elimination or direct DMA paths. We cover (i) architecture-level control-plane evolution, (ii) data-path mechanisms from local to remote storage, (iii) deployment models including disaggregated and fabric-attached resources, and (iv) workload-driven requirements from AI training pipelines and checkpoint-heavy execution. This broader framing connects transport/protocol building blocks (NVMe, NVMe-oF, RDMA) with filesystem/runtime behavior and application-level bottlenecks in one system narrative \cite{nvmeof_spec,pcie_spec,liu2017rackscale,murray2021tfdata}.

Our contributions are fourfold:
\begin{enumerate}
    \item \textbf{Systematic co-evolution view.} We present a structured account of how storage and GPU subsystems have co-evolved from CPU-orchestrated paths to accelerator-participatory designs \cite{nvidia_gds_blog,silberstein2013gpufs,qureshi2023bam}.
    \item \textbf{Unified taxonomy.} We develop a taxonomy of GPU-integrated I/O architectures spanning control-plane ownership, data-path realization, deployment context, and consistency/metadata semantics.
    \item \textbf{Cross-system design synthesis.} We synthesize recurring principles and trade-offs across research prototypes, production frameworks, and standards-backed deployments \cite{geminifs2025,nvidia_gds_overview,infiniband_spec}.
    \item \textbf{GPU-driven paradigm outlook.} We analyze emerging directions in which GPUs become active storage participants, and we identify open problems for end-to-end system redesign in AI/HPC environments \cite{rajbhandari2021zeroinfinity,zhang2025hpcstoragesurvey}.
\end{enumerate}
